\chapter{Globes and Maps}
\section{\textbf{Get Rolling}}
\textbf{Locate India and five other countries below.}
\begin{enumerate}
    \item Indonesia
    \item China
    \item Japan
    \item Sri Lanka
    \item Pakistan
\end{enumerate}
\section{\textcolor{red}{What to Know}}
Globes and maps are invaluable tools that help us to explore and understand the world around us.
\section{\textbf{Learning Objectives}}
\begin{itemize}
    \item to describe the shape of Earth
    \item to name and locate different continents and oceans on a world map
    \item to explain different types of maps
    \item to acquire the skills of reading map
    \item to analyse the advantages and disadvantages of using maps and globes
\end{itemize}
\section{I Explore}
Hundreds of years ago, people believed that Earth was flat and that the \textcolor{red}{horizon} was the edge or end of Earth. They thought if sailors reached the edge, they could fall off the edge.
\vspace{0.1in}
\par
Then how did we come to know that Earth was spherical?
\vspace{0.1in}
\par
Ferdinand Magellan, a Portuguese explorer, became the first European to sail around the world and proved that Earth was round. A Greek thinker, Pltheato, who lived almost 2450 years ago, suggested that Earth was spherical. Later almost 2374 years ago, Aristotle, a student of Plato, declared that Earth was a sphere.
\vspace{0.6in}
\par
The picture of Earth taken from space in recent times also proved beyond doubt that it is spherical. It is round but slightly flattened at the top and the bottom.
\vspace{0.1in}
\par
\textcolor{red}{I wonder} how humans figured out the correct shape of Earth before the pictures of the planet could be taken from space!
\vspace{0.6in}
\section{\textcolor{red}{Globe}}
Earth is a big planet. It is not possible to see the entire Earth from one place. To study Earth, we need its model. A \textbf{globe} is a small 3D Model of earth. A globe also helps us to know the true shape of the planet without to going space.
\vspace{0.1in}
\par
A globe shows how Earth rotates on an imaginary line passing through its centere. This line is called the \textbf{axis.} The points through which the axis of Earth passes are the North Pole and the South Pole. The North Pole is the northernmeost point of Earth and the South Pole is southernmeost point of Earth. Earth is tilted and rotates on its own axis from west to east.
\vspace{0.4in}
\par
At a time, we can see only one half of a globe. To see its other half, of a globe around on its axis.
\vspace{0.1in}
\par
A globe shows the relative sizes and locations of seas, oceans, continents and countries on it. Therefore, a globe shows the distribution of land and water on the earth’s surface.
\vspace{0.1in}
\par
A globe is ineed the most \textcolor{red}{accurate} way of looking at the sizes, shapes and locations of places on Earth. However, it have some disadvantages.
\begin{itemize}
    \item A globe is bulky. Therefore, it is difficult to carry it around.
    \item It shows only one side of Earth at one time.
    \item The details given are very small and not clear. Therefore, it is not helpful if we were to study part of Earth, say a country, in detail.
    \item It can be neither rolled nor carried in a small bag.
\end{itemize}
\section{\textcolor{red}{Continents and Oceans}}
\textbf{Continents} are large landmasses on Earth. There are seven continents in the world. Most continent have many countries. Let us find our more about these continents.
\begin{itemize}
    \item \textbf{Asia} is the biggest continent in the world. It covers one-third of the planet’s total landmasses. Almost 60 percent of the world's population lives here. India is the part of Asia. Mount Everest, the hightst mountain peak in the world, is located in Asia.
    \item \textbf{Africa} is the second largest continent in the world. It has the largest number of countries in it. The Sahara Desert, the world's largest desert, and the Nile River, the world's longest river, are situated in Africa.
    \item \textbf{North America} is the third largest continent. The world's largest island, Greenland, is situated in the northerneast of North America.
    \item \textbf{South America} is the fourth largest continent. The Amazon rainforests in this continent are called the lungs of Earth. The Amazon River is in the northern part of South America.
    \item \textbf{Antarctica} is the southern part of Earth. It encircles the South Pole. Humans do not live here, as the continent is permanently covered with thick layers of ice. Antarctica is also called the white continent. Some scientists stay there temporarily to study more about it.
    \item \textbf{Australia} is the smallest continent in the world. It is also known as the 'island continent' as it is surrounded by water on all sides. It is the only continent named after a country-Australia.
    \item \textbf{Europe} is the second smallest in the world. The world's smallest country, the Vatican City, is in Europe. The URL Mountain range separates it from Asia.
\end{itemize}
The continents are separated from one another by large water bodies are called \textbf{oceans.} Oceans are very deep water bodies. More that half of the earth's surface is covered with oceans.
\vspace{0.1in}
\par
There are five oceans in the world. Let us find more about these oceans.
\begin{itemize}
    \item The \textbf{Pacific Ocean} is the largest and deepest ocean in the world.
    \item The \textbf{Antlantic Ocean} is the second largest and busiest ocean in the world. It is shaped like an 'S'.
    \item The \textbf{Indian Ocean} is the third largest ocean in the world. The Arabian Sea and the Bay of Bengal are parts of the Indian ocean.
    \item The \textbf{Antarctic Ocean} is the ocean that encircles Antarctica. It is also known as the Southern Ocean. It is the coldest ocean in the world.
    \item The ocean around the North Pole is the \textbf{Arctic Ocean.} This is the smallest ocean in the world and remains frozen in winters.
\end{itemize}
\vspace{0.26in}
\section{\textcolor{red}{Maps}}
A globe can study Earth. But when we want to study a country or a state, a globe isn't of much help. In such a situation, people prefer using maps. A \textbf{map} is a representation or drawing of the earth's surface or a part of it drawn on a flat surface. Small areas can be represented accurately on a map.
\vspace{0.1in}
\par
Maps shows places and distances on a flat surface or paper. A large area of land can be shown on a small map. A map provides imformation about human-made features as well as natrual features. A book of maps is called an \textbf{atlas.} It is a collection of different types of maps.
\vspace{0.1in}
\par
Maps are very easy to use. However, it is difficult to show the spherical Earth on a flat piece of paper. If we try to peel off the whole skin of an orange and flatten it, it will many wrinkles and cuts. Similarly, if we flatten Earth for a map projection, distortions in terms of shape, distance, direction or land area are bound to happen.
\subsection{\textcolor{orange}{Types of Maps}}
Maps are of different types of according to the imformation they give; for example, political maps, road maps, physical maps and thematic map etc. Maps use different colours and symbols to show different things on the surface of the earth.
\subsection{\textbf{Road map}}
A road map shows national and state highways. Road maps are useful for drivers to find their way around. Road maps depicts how major roads connect with various cities and towns.
\vspace{0.1in}
\par
Road maps are not used much these days. Modern technology provides the \textbf{GPS} or \textbf{Global Positioning System.} It is a network of satellites that give accurate locations and the time required to reach our \textcolor{red}{destination}.
\subsection{\textbf{Political map}}
A political map shows boundaries that define the territories of countries, states, cities, towns and villages.
\subsection{\textbf{Thematic map}}
A thematic map focuses on one theme or topic such as rainfall, population and vegetation.
\subsection{\textbf{Physical map}}
A physical map shows the natural or physical features of the surface of the earth like mountains, plains, plateaus, deserts, and water bodies. Generally, these maps are very colourful, as they show different features in different colours for easy reading.
\subsection{\textcolor{orange}{Elements of a Map}}
We must know how to read a map in order to understand the information it provides. Every map has some essential elements like direction indicators, map scale, symbols and keys. These elements make it easy to read a map.
\subsection{\textbf{Directions}}
There are four major directions-North, South, East and West. They are known as \textbf{cardinal points.} You must have seen that most maps have an arrow marked with the letter 'N' at the top right-hand corner. It is the \textbf{North pointer.} It tells us the north direction. You can easily find other directions form this reference point. The bottom of the map is south. To the left is the west and to the right is the east. There are four \textbf{intermediate directions.}
\begin{itemize}
    \item North-East (NE) lies between North and East.
    \item South-East (SE) lies between South and East.
    \item South-West (SW) lies between South and West.
    \item North-West (NW) lies between North and West.
\end{itemize}
It is possible to locate any place accurately with the help of the intermediate directions. The compass is a \textcolor{red}{device} that indicates direction. It is a very important instrument used in navigation. It has a magnetic needle that always points north.
\subsection{\textbf{Scale}}
A map is much smaller than the actual size of the place it represents. To ensure the distance accurately on a map, a map scale is used.
\vspace{0.1in}
\par
The scale of a map is a ratio between the actual distance on the map and the actual distance on the ground. Suppose the distance between your house and your grandmother's house is 100 km. It is not possible to show the exact distance on a map. You can easily represent this distance on the map by a smaller unit like 10 cm. In such a case, the scale will be as follows:
\vspace{0.1in}
\par
1 km = 10 km
\vspace{0.1in}
\par
The scale on the map is shown as a line with distance marked on it or as 1 cm : 10 km.
\vspace{0.1in}
\par
A small-scale map can show large areas like forests or countries. A large-scale map shows small areas such as a village, a town, a city or a neighbourhood.
\subsection{\textbf{Symbols}}
Different things like railway lines, schools, hospitals, forests and more can be shown on a map. But maps have limited space. Then, how can we show these things on a map? We can do this using symbols. A symbol is a representation of an object on a map. These symbols making easy and enjoyable. To show a railroute on a map, we can use broken lines, in the form of dots. Similarly, we can use a tree symbol to show a forest on a map.
\subsection{\textbf{Colours}}
Various colours are used to represent different features on a map. For example, blue is used to represent water bodies, green to represent plains, yellow to represent pleteaus and brown to represent mountains. A key is given on the map to help readers understand what different symbols mean. A key is also known as a legend.